\begin{table}[htbp]
\centering
\begin{threeparttable}
\caption{Alternative Policy Variants (APV) — counterfactual price effects}
\label{tab:v3_apv}
\small
\begin{tabular}{lrrrrr}
\toprule
Class & $\bar p_{S_1}$ & V0 (v3) & V1 (50\%) & V2 (no entry) & V3 (10\% pref.) \\
\midrule
non-pharma & 0.759 & +0.2585 & +0.0925 [35.8\%] & +0.3930 [152.0\%] & -0.0036 [-1.4\%] \\
pharma & 0.656 & +0.3079 & +0.1032 [33.5\%] & +0.4847 [157.4\%] & +0.0015 [0.5\%] \\
\bottomrule
\end{tabular}
\begin{tablenotes}\footnotesize
\item Price $\Delta$ relative to observed $S_1$ (open, Pre pool).
V0 = full 100\% SME set-aside with endogenous Post pool (main text).
V1 = partial 50/50 mix. V2 = SME-only but Pre pool (isolates intensive
margin). V3 = open auction with 10\% price preference for SMEs.
Bracketed values are each variant's $\Delta$ as \% of V0's $\Delta$,
so V1 = 50\% means ``half the v3 effect''.
\end{tablenotes}
\end{threeparttable}
\end{table}
